\begin{table}[t]
\centering
\begin{threeparttable}
\caption{Letting the macro volatility regime load differentially on Bitcoin}
\label{tab:macro}
\begin{tabular}{lc}
\toprule
 & Fisher-$z$ correlation with SPY \\
\midrule
$\hat\tau$: BTC $\times$ Post & 0.0195 \\
 & (0.0201) \\
VIX $\times$ BTC & 0.0100*** \\
 & (0.0018) \\
MOVE $\times$ BTC & -0.00034 \\
 & (0.00029) \\
Randomization $p$ on $\hat\tau$ (placebo-in-space) & 0.80 \\
SD of the randomization distribution & 0.0536 \\
Within $R^2$ & 0.0083 \\
$N$ & 540 \\
\bottomrule
\end{tabular}
\begin{tablenotes}[flushleft]
\footnotesize
\item The specification adds the monthly VIX and MOVE levels interacted with the Bitcoin indicator to the baseline coin-month regression. The main effects are collinear with the calendar-month fixed effects and are absorbed, so what these interactions identify is whether the macro volatility regime loads differentially on Bitcoin relative to the never-listed control coins. It does: Bitcoin's equity co-movement is a strong positive function of the VIX level, while the rate-volatility interaction does nothing. Including these interactions moves the treatment coefficient from $-0.0212$ in the baseline to the value shown here, a shift of roughly two-thirds of a randomization standard deviation, and reverses its sign. Both estimates lie deep inside the randomization distribution, so neither is distinguishable from zero and the sign reversal is not a finding; it is a statement about precision, and it is the reason the paper describes the estimate as robustly null rather than robustly negative.
\end{tablenotes}
\end{threeparttable}
\end{table}
